% ===== SECTION 3: BARRIERS TO ENTRY =====
\section{\texorpdfstring{\color{MidnightBlue}Inquiries Regarding Potential Barriers to Entry}{Barriers to Entry}}

\begin{sidebar}{RFI Reference}
Section C: Questions to Field --- Inquiries Regarding Potential Barriers to Entry (Questions 1--4)
\end{sidebar}

% ----- Question 1 -----
\subsection{Question 1}
\textit{``What, if anything, may prevent your company from bidding on this scope of work?''}

\placeholder{Address Sentry's specific barriers honestly. Key points to cover:}

\begin{itemize}
    \item \placeholder{\textbf{Fleet \& Drivers:} Sentry does not own or operate a bus fleet and does not directly employ CDL drivers. Our participation depends on whether RIDE structures the RFP to include a separate technology/system management component, or whether we can partner with a bus operator.}
    \item \placeholder{\textbf{RI Presence:} Sentry does not currently have an operational footprint in Rhode Island. If yard space, fuel access, or local infrastructure are preconditions for bidding, state-provided support (per Question 4) would be critical.}
    \item \placeholder{\textbf{Procurement Structure:} If the RFP is structured exclusively for vertically integrated bus operators (i.e., one vendor provides buses, drivers, AND technology), Sentry would be unable to bid independently. A structure that separates technology/system management from vehicle operations would enable our participation.}
\end{itemize}

% ----- Question 2 -----
\subsection{Question 2}
\textit{``What contract structure(s) would make participation in RIDE's Statewide Transportation System (STS) more feasible for your organization? For example, would any of the following potential contract structures make it more likely that your organization would bid on this scope of work?''}

\vspace{0.2cm}

Sentry strongly encourages RIDE to consider the following contract structures, which would increase competition, attract specialized providers, and ultimately deliver better outcomes for Rhode Island students:

\begin{description}[style=nextline, leftmargin=0.5cm]
    \item[\faIcon{check-circle} \textbf{Separate Technology / System Management Contract:}] \placeholder{Advocate for a standalone technology procurement---a ``Statewide System Manager'' contract covering routing optimization, dispatch, fleet analytics, video monitoring, and parent communication. This allows RIDE to select best-in-class technology independently of bus operator selection, ensures vendor-neutral oversight across all zones, and avoids the conflict of interest when operators self-report performance data.}

    \item[\faIcon{check-circle} \textbf{Multiple Vendor Awards per Zone:}] Allowing multiple vendors per zone would enable technology providers like Sentry to partner with bus operators, bringing specialized routing and monitoring capabilities to zones where operators may lack in-house technology.

    \item[\faIcon{check-circle} \textbf{Allowing Subcontracting or Partnerships:}] Cross-zone partnerships would enable smaller or specialized firms to combine resources. Sentry could serve as the technology backbone for one or more bus operators across multiple zones, creating a more competitive field.

    \item[\faIcon{check-circle} \textbf{Subdividing Larger Zones:}] Smaller zones lower the fleet threshold for entry, encouraging regional and minority-owned operators to participate. Sentry's routing platform can optimize routes regardless of zone boundaries, ensuring efficiency is not lost through subdivision.
\end{description}

\begin{sidebar}{Recommendation to RIDE}
RIDE already operates with a separate System Manager---Transpar currently provides routing, program billing, vendor performance monitoring, and customer service support across all 49~LEAs. The RFP is an opportunity to \textit{strengthen} this model by procuring an upgraded technology platform with modern capabilities: AI-powered routing optimization, real-time GPS dispatch, live video violation detection, and automated cost analytics. A next-generation System Manager would replace the current 700-check-per-year manual billing process, eliminate annual paper re-enrollment that contributes to the 550 ``ghost rider'' problem, and give RIDE independent performance data rather than relying on operator self-reports.
\end{sidebar}

% ----- Question 3 -----
\subsection{Question 3}
\textit{``What minimum route density or fleet size makes bidding on a zone more viable for your organization?''}

\placeholder{As a technology provider, Sentry's viability is not tied to route density or fleet size in the traditional sense. Our platform scales from 29 routes (Zone D) to 122+ routes (Zone A1) with equal efficiency. Address this from both angles:}

\begin{itemize}
    \item \placeholder{\textbf{As technology provider:} No minimum --- the platform is zone-agnostic and scales linearly. In fact, deploying across ALL zones is most efficient, as it enables cross-zone optimization.}
    \item \placeholder{\textbf{If partnering with a bus operator:} The partner operator's fleet and driver minimums would apply. Generally, 25--30 routes represents a viable minimum for a dedicated depot.}
\end{itemize}

% ----- Question 4 -----
\subsection{Question 4}
\textit{``What support (e.g., yard space, fuel access, maintenance facilities, etc.) could the State provide or assist with that would make it more likely for your organization to bid?''}

\begin{description}[style=nextline, leftmargin=0.5cm]
    \item[\faIcon{database} \textbf{Data Access:}] \placeholder{Providing prospective bidders with anonymized route data, student origin/destination data (FERPA-compliant), current schedules, and historical performance metrics would allow technology providers to demonstrate value through optimized routing scenarios during the RFP process.}

    \item[\faIcon{server} \textbf{Technology Infrastructure:}] \placeholder{State-provided or state-facilitated access to API integrations with existing RIDE systems (student information systems, IEP databases) would accelerate deployment and reduce implementation risk.}

    \item[\faIcon{warehouse} \textbf{Yard Space \& Facilities:}] \placeholder{State-owned or state-leased yard space available to new entrant operators would reduce the capital barrier for bus operators partnering with technology providers like Sentry.}

    \item[\faIcon{gas-pump} \textbf{Fuel Access:}] \placeholder{Centralized fueling contracts or bulk purchasing arrangements could reduce costs for smaller operators and level the playing field against incumbents with established fueling infrastructure.}

    \item[\faIcon{clock} \textbf{Extended Transition Period:}] \placeholder{A transition period longer than 8 months (November 2026 award to July 2027 service start) would give new entrants time to establish local operations, recruit drivers, and deploy technology.}
\end{description}
